\section{Stable tubes and a kernel class}\label{sec:tubes}

The output of this section is a nonzero family of smooth tube data over
$S^9$ whose associated stable spherical fibration is trivial. We first
specify the function model and its stabilizations, then obtain the class
from the stable smooth $h$-cobordism space.

\subsection{Functions and domains}

For positive integers $k,l$, write
\[
 q_{k,l}(x,y)=-\norm{x}^{2}+\norm{y}^{2},
 \qquad (x,y)\in\R^k\oplus\R^l.
\]
Let $T_{k,l}$ be the path component of $q_{k,l}|_{S^{k+l-1}}$ in
\[
 \{h\in C^1(S^{k+l-1},\R):0\text{ is a regular value of }h\},
\]
with its $C^1$ topology. We also regard its elements as homogeneous functions
\begin{equation}\label{eq:homogeneous-extension}
 g(ru)=r^2h(u),\qquad r>0,\quad u\in S^{k+l-1},\qquad g(0)=0.
\end{equation}
Indeed $g$ is $C^1$ on the whole vector space, since $g(w)=O(\norm{w}^2)$
and $dg(w)=O(\norm{w})$, and
\begin{equation}\label{eq:angular-gradient}
 \nabla g(ru)=r\bigl(2h(u)u+\nabla_Sh(u)\bigr).
\end{equation}
The radial and tangential terms are orthogonal. Consequently $0$ is a
regular value of $h$ precisely when $g$ has no critical point away from $0$.
When $h$ is smooth, $g$ is smooth off the origin, and its negative domain
\[
 D(h)=\{u\in S^{k+l-1}:h(u)\leq0\}
\]
is a smooth tube: a domain isotopic to
$D(q_{k,l})\cong S^{k-1}\times D^l$.

We will pass between functions and domains in families. Here are the
details of that passage. For a fixed smooth domain $D$ with cooriented
boundary, the smooth functions defining $D$ have negative sign in
$\Int D$, positive sign outside $D$, and positive outward normal derivative
on $\partial D$. Their convex combinations retain all three properties.
Thus the space of defining functions for $D$ is convex. A collar of
$\partial D$ supplies one such function. For nearby domains, express the
boundaries as normal graphs and extend their graph isotopies to ambient
isotopies supported in the collar. Pulling back the chosen function gives
local continuous choices of defining functions. Over a compact parameter
space a partition of unity combines these choices: at each parameter all
the terms define the same domain, so convexity applies. This construction
also works relative to a subspace on which defining functions are already
chosen, by interpolation with that choice in a neighborhood. In particular,
a homotopy of smooth domains lifts to a homotopy of defining functions,
with prescribed endpoints.

Smooth angular functions and $C^1$ angular functions give the same
homotopy groups. To see this directly, for a compact family put
\[
 \mu=\min_{b,u}\bigl(4h_b(u)^2+\norm{\nabla_Sh_b(u)}^2\bigr)^{1/2}>0.
\]
Sufficiently close approximation in angular $C^1$ norm preserves this
inequality, also throughout straight interpolation. Smooth finitely many
angular functions and take a smooth partition of unity in the parameters
to obtain such an approximation. Cutoffs in the parameters give the
relative version. Applying the same argument to a homotopy proves both
surjectivity and injectivity on homotopy groups. These descriptions agree
with the smooth tube models of \cite[\S3.6]{AIS2026}; see also
\cite[Definitions~3.10 and~3.36]{AACK2025} for homogeneous $C^1$
functions and the tube-type space with its $C^1$ topology.

Direct sum is well defined in this model:
\[
 (g\oplus g')(w,z)=g(w)+g'(z).
\]
Its gradient vanishes only when both block gradients vanish, hence only at
$(w,z)=0$. Direct sums of paths from the standard forms also show that the
indices add. Although $g\oplus g'$ can fail to be smooth on a coordinate
axis, it is $C^1$ there; subsequent angular smoothing preserves its tube
class. Adding positive squares has the usual geometric interpretation.
In the sphere chart $u=w/\norm{w}$, $v=z/\norm{w}$, its negative domain is
\[
 h(u)+\norm{v}^2\leq0.
\]
This is the disk thickening of the old negative domain, with the boundary
rounded. Adding negative squares gives the corresponding construction on
the positive domain.

Fix the coordinate identifications putting negative coordinates before
positive ones. Additions of positive and negative squares then form a
commuting two-parameter stabilization system. The diagonal is cofinal, and
we use the based mapping telescope
\begin{equation}\label{eq:stable-tubes}
 \T=\operatorname*{hocolim}_{m}
 \bigl(T_{m,m}\xrightarrow{\,\oplus q_{1,1}\,}T_{m+1,m+1}\bigr).
\end{equation}
In particular, a nullhomotopy after adding any fixed numbers of positive
and negative squares implies stable nullity: adding the missing squares
reaches a diagonal stage.

\subsection{The stable kernel}

Let $G$ be the stable space of homotopy self-equivalences of spheres,
stabilized by suspension; $BG$ classifies stable spherical fibrations.
The negative domains of a tube family have the homotopy type of spheres,
and their stable fiber homotopy class defines a map $\gamma:\T\to BG$.
To describe its fiber, let $H(X)$ be the smooth $h$-cobordism space
of a compact smooth manifold $X$, with identified incoming end and
prescribed product vertical boundary. Its stable version is
\[
 \Hs(X)=\operatorname*{hocolim}_{r}H(X\times J^r),
\]
where $J$ is a closed interval and the maps are interval product
stabilizations. For $X=*$ this is equivalently the stabilization of
disks of increasing dimension. We will specify the collared product
model when applying it in Section~\ref{sec:obstruction}.

We use the stable parametrized $h$-cobordism theorem
\cite[Theorems~0.1 and~0.3]{WJR2013} in the form
\begin{equation}\label{eq:stable-hcob-theorem}
 \Hs(X)\simeq\Omega\Wh^{\mathrm{diff}}(X).
\end{equation}
Here
$\Wh^{\mathrm{diff}}(X)$ denotes the infinite loop space of the cofiber
of the unit $\Sigma^\infty X_+\to A(X)$, where $A(X)$ is Waldhausen's
algebraic $K$-theory spectrum of $X$.

Waldhausen's tube theorem identifies the fiber of $\gamma$ with
$\Hs(*)$.

\begin{theorem}[Waldhausen's tube fibration]\label{thm:tube-fibration}
For the stable smooth tube space \eqref{eq:stable-tubes} there is a
homotopy fiber sequence
\begin{equation}\label{eq:tube-fibration}
 \Hs(*)\longrightarrow\T\xrightarrow{\gamma}BG.
\end{equation}
The rigid tube map $\iota:BO\to\T$, defined by nondegenerate quadratic
forms, satisfies $\gamma\circ\iota\simeq BJ$, where
$BJ:BO\to BG$ is the stable spherical fibration map.
\end{theorem}

This is the smooth tube theorem of \cite{Waldhausen1982}, with the
function and domain comparison and the rigid map described in
\cite[\S3.6, Theorem~3.30 and Remark~3.34]{AIS2026}. The negative-domain
frame model in \cite[\S4.1, Definition~4.3, Propositions~4.4, 4.6, 4.7 and Remark~4.5]{AACK2025}
identifies this particular map $\gamma$ with the classifying map of the
negative-domain spherical fibration. To identify the rigid composite,
retract the negative domain of a quadratic form onto
the unit sphere of its negative eigenspace. Stabilizing the corresponding
Grassmannians gives $BO$, and taking those sphere bundles is exactly
$BJ$. The comparison of the partition and spherical-domain models
is a stable equivalence; we use no finite-stage equivalence between
those models. All equivalences and fiber sequences here are understood
up to weak homotopy equivalence.

\begin{proposition}\label{prop:kernel-class}
There exists
\begin{equation}\label{eq:kernel-class}
 0\ne a\in\pi_9\T,\qquad \gamma_*a=0.
\end{equation}
Moreover, $(BJ)_*:\pi_9BO\to\pi_9BG$ is injective.
\end{proposition}

\begin{proof}
Bott periodicity gives $\pi_8O\cong\pi_9O\cong\Z/2$
\cite[Corollary to Theorem~II, equation~(1.5)]{Bott1959}.
The stable $J$-homomorphisms in both degrees are injective
\cite[Theorem~9.5]{Adams1966}. We also use
\[
 \pi_9^S\cong(\Z/2)^3,
 \qquad
 \bigl(\pi_{10}\Wh^{\mathrm{diff}}(*)\bigr)_{(2)}
 \cong\Z/8\oplus(\Z/2)^2.
\]
The first value appears in \cite[Table~1, degree~9]{IWX2023}; the second
is Rognes's computation \cite[Theorem~5.8]{Rognes2002}, also recorded in
\cite[Theorem~3.3]{Rognes2014}. Only the two-primary assertion is used.

In positive degrees $\pi_iG=\pi_i^S$. One may use either based or
unbased sphere equivalences here, since the evaluation fibration has
base a sphere whose connectivity tends to infinity under stabilization.
The relevant part of the long exact sequence of
\eqref{eq:tube-fibration} is
\begin{equation}\label{eq:tube-long-exact}
 \pi_{10}\T\longrightarrow\pi_{10}BG
 \xrightarrow{\partial}\pi_9\Hs(*)
 \longrightarrow\pi_9\T\xrightarrow{\gamma_*}\pi_9BG.
\end{equation}
The image of $\pi_{10}BO$ in $\pi_{10}BG=\pi_9^S$ is the
order-two image of $J_9$. Theorem~\ref{thm:tube-fibration} puts it in
$\ker\partial$. Since $\pi_9^S$ has order eight,
$\abs{\im\partial}\leq4$. On the other hand,
\eqref{eq:stable-hcob-theorem} identifies
\[
 \pi_9\Hs(*)=\pi_{10}\Wh^{\mathrm{diff}}(*),
\]
whose two-primary part has order $32$. Thus $\partial$ is not surjective.
Exactness in \eqref{eq:tube-long-exact} implies that the image in
$\pi_9\T$ is nonzero and is contained in $\ker\gamma_*$, proving
\eqref{eq:kernel-class}. Finally, $\pi_9BO\to\pi_9BG$ is $J_8$, which
is injective by the other degree of Adams's theorem cited above.
\end{proof}

\begin{lemma}\label{lem:localized-family}
Fix $B=S^9$ and a standard height function $\theta:B\to\R$. The class
$a$ in \eqref{eq:kernel-class} has, for arbitrarily large $m\geq20$,
a representative $b\mapsto h_b$ in $T_{m,m}$ such that
$(b,u)\mapsto h_b(u)$ is smooth and $h_b=q_{m,m}|_{S^{2m-1}}$ on
neighborhoods of both critical points of $\theta$. Its homogeneous
extension $g_b$ satisfies, for constants $0<c\leq C<\infty$,
\begin{equation}\label{eq:uniform-homogeneous-gradient}
 c\norm{w}\leq\norm{\nabla_wg_b(w)}\leq C\norm{w}
 \qquad(b\in B,\ w\in\R^{2m}).
\end{equation}
For any fixed metric on $B$, there is also a constant $C_1$ such that
\begin{equation}\label{eq:base-mixed-bound}
 \norm{d_wd_bg_b(w)}\leq C_1\norm{w}.
\end{equation}
The derivatives in \eqref{eq:base-mixed-bound} extend continuously to
$w=0$ with value zero.
\end{lemma}

\begin{proof}
A map from the compact sphere to the mapping telescope
\eqref{eq:stable-tubes} has bounded height. A finite initial portion of
the telescope retracts onto its last stage. Hence $a$ is represented by
a based map $S^9\to T_{m,m}$ at a finite stage. Further stabilization
preserves its stable class and makes $m$ arbitrarily large.

Choose an oriented disk in $B$ whose closure avoids the two critical
points of $\theta$. Collapsing the complement of its interior gives a
degree-one map $B\to S^9$, constant at the target basepoint off that
disk. Equivalently one can take a smooth map of degree one supported
there, obtained by a relative smooth approximation of this collapse
map. Precomposition preserves the based homotopy class, since a
degree-one self-map of $S^9$ is based homotopic to the identity. The
resulting family is standard on neighborhoods of both critical points.

Apply the relative angular approximation described above, retaining the
standard function on smaller such neighborhoods. More explicitly,
approximate the continuous map into $C^1(S^{2m-1})$ by a finite smooth
partition of unity in $b$ multiplying smooth angular functions; then
multiply its difference from $q_{m,m}$ by a parameter cutoff which is
zero on the smaller neighborhoods and one outside the original
standard neighborhoods. The approximation remains uniformly close in
angular $C^1$ norm. The positive regularity margin and straight
interpolation show that its stable class is still $a$ and its spherical
class is zero.

For this smooth family, the continuous positive function
$\norm{2h_b(u)u+\nabla_Sh_b(u)}$ has a positive minimum and a finite
maximum on $B\times S^{2m-1}$. Formula~\eqref{eq:angular-gradient}
gives \eqref{eq:uniform-homogeneous-gradient}, including the origin.
Finally $d_bg_b(ru)=r^2d_bh_b(u)$. Differentiating in $w$ gives a
degree-one expression whose angular coefficients are bounded by
compactness, proving \eqref{eq:base-mixed-bound}. The degree-two value
bound and degree-one derivative bound also prove the asserted $C^1$
extension at zero. As before, the nonzero based stable class cannot
be freely nullhomotopic, because the fundamental group action fixes
zero.
\end{proof}
